\section{Image preparation and mathematical encoding model}

\subsection{Coordinate system and framed source}

Let the loaded source image be defined on pixel coordinates $(u,v)$. Percepta first constructs an output canvas of width $W$ and height $H$. A point $(x,y)$ on that canvas is mapped back to the source through crop, zoom, translation and rotation.

Using homogeneous coordinates, the combined inverse mapping can be written as

\begin{equation}
\begin{bmatrix}
u\\v\\1\end{bmatrix}
=
\mathbf{T}_{\mathrm{source}}
\mathbf{R}(-\varphi)
\mathbf{S}\!\left(\frac{1}{z}\right)
\mathbf{T}_{\mathrm{pan}}
\begin{bmatrix}x\\y\\1\end{bmatrix},
\label{eq:framing}
\end{equation}

where $z$ is the zoom factor, $\varphi$ is the rotation angle, and the translation matrices include crop centring and the user-defined horizontal and vertical pan. The source value at non-integer $(u,v)$ is obtained by interpolation. Bilinear interpolation is sufficient for interactive work; higher-order resampling may be used for final export.

The framed RGB image is denoted

\begin{equation}
\mathbf{C}(x,y)=
\begin{bmatrix}
R(x,y)\\[1mm]G(x,y)\\[1mm]B(x,y)
\end{bmatrix},
\qquad 0\le R,G,B\le 1.
\end{equation}

Values are normalised to the interval $[0,1]$ before rendering.

\subsection{Contrast preprocessing}

The \control{Source contrast} parameter changes the distance of each channel value from a mid-grey pivot. A convenient operational form is

\begin{equation}
C'_q(x,y)=\clamp\!\left[
\frac{1}{2}+\gamma_c\left(C_q(x,y)-\frac{1}{2}\right),\,0,\,1
\right],
\qquad q\in\{R,G,B\},
\label{eq:contrast}
\end{equation}

where $\gamma_c=1$ leaves contrast unchanged, $\gamma_c>1$ increases contrast and $0<\gamma_c<1$ compresses it. The function

\begin{equation}
\clamp(a,0,1)=
\begin{cases}
0,&a<0,\\
a,&0\le a\le 1,\\
1,&a>1
\end{cases}
\end{equation}

prevents invalid channel values. An implementation can use an equivalent image-library contrast operator, but Equation~\eqref{eq:contrast} captures the intended behaviour.

Contrast is applied before geometric encoding. It therefore changes the distribution of widths or radii, not merely the appearance of a completed result. High values expand the number of locations that map to minimum or maximum geometry and can remove subtle tonal information.

\subsection{Channel-specific geometric coverage}

For every channel $q$, Percepta creates a geometric mask

\begin{equation}
M_q(x,y)\in[0,1].
\end{equation}

A value of zero means that channel contributes no light at $(x,y)$; a value of one means full channel coverage. Antialiased edges take intermediate values.

The source channel is converted into a modulation amplitude by

\begin{equation}
a_q(x,y)=\sat\!\left(a_{\min}+s\,C'_q(x,y)\right),
\label{eq:amplitude}
\end{equation}

where $s$ is related to \control{Strength}, $a_{\min}$ is the renderer's minimum visible geometry and

\begin{equation}
\sat(a)=\clamp(a,0,1).
\end{equation}

In practice, the geometry may use a nonlinear response

\begin{equation}
a_q(x,y)=\sat\!\left(a_{\min}+s\,[C'_q(x,y)]^{\eta}\right),
\label{eq:gamma}
\end{equation}

where $\eta$ controls how quickly small channel values become visible. The current controls expose Strength rather than $\eta$ directly; the equation is useful when comparing renderer implementations or future calibration modes.

\subsection{RGB displacement}

Colour separation is produced by evaluating or placing the channel structures at slightly different positions. Let the separation magnitude in output pixels be $d$. A symmetric three-channel arrangement can be represented by displacement vectors

\begin{equation}
\boldsymbol{\delta}_R=d\begin{bmatrix}-1\\0\end{bmatrix},\qquad
\boldsymbol{\delta}_G=d\begin{bmatrix}0\\0\end{bmatrix},\qquad
\boldsymbol{\delta}_B=d\begin{bmatrix}1\\0\end{bmatrix},
\label{eq:rgbshift}
\end{equation}

or by the same vectors rotated into the local normal direction of a stripe or spiral. The renderer then uses

\begin{equation}
\widetilde{M}_q(x,y)=M_q\!\left(x-\delta_{q,x},y-\delta_{q,y}\right).
\end{equation}

The exact orientation can be pattern-dependent. What matters perceptually is that channel supports no longer coincide perfectly. Where two channels overlap, additive secondary colours appear: red plus green gives yellow, green plus blue gives cyan, and red plus blue gives magenta.

\subsection{Additive composition}

The final RGB output is the vector image

\begin{equation}
\mathbf{I}_{\mathrm{out}}(x,y)=
\begin{bmatrix}
\widetilde{M}_R(x,y)\\[1mm]
\widetilde{M}_G(x,y)\\[1mm]
\widetilde{M}_B(x,y)
\end{bmatrix}.
\label{eq:additive}
\end{equation}

On a black background, each mask directly drives the corresponding display primary. If an opacity or background term is introduced, a more general compositing expression is

\begin{equation}
\mathbf{I}_{\mathrm{out}}=
\boldsymbol{\alpha}\odot\mathbf{M}
+\left(\mathbf{1}-\boldsymbol{\alpha}\right)\odot\mathbf{I}_{\mathrm{bg}},
\end{equation}

where $\odot$ denotes element-wise multiplication. The standard Percepta appearance uses a dark background because it maximises the visual distinction between individual RGB elements while preserving additive mixing.

\subsection{Area-average preservation}

Consider a local pattern cell $\Omega$ whose dimensions are small compared with the scale of source variation. The mean output colour is

\begin{equation}
\overline{\mathbf{I}}_{\Omega}
=
\frac{1}{|\Omega|}
\iint_{\Omega}
\mathbf{I}_{\mathrm{out}}(x,y)\,\mathrm{d}x\,\mathrm{d}y.
\label{eq:localmean}
\end{equation}

A well-calibrated renderer chooses stripe widths or dot areas so that

\begin{equation}
\overline{\mathbf{I}}_{\Omega}
\approx
\begin{bmatrix}
\overline{R}_{\Omega}\\[1mm]
\overline{G}_{\Omega}\\[1mm]
\overline{B}_{\Omega}
\end{bmatrix}.
\label{eq:preserve}
\end{equation}

This is the central conservation principle. The pattern can differ radically from the source at individual pixels while preserving average channel energy over a neighbourhood.

\begin{warningbox}
Exact equality in Equation~\eqref{eq:preserve} is not guaranteed everywhere. Minimum widths, clipping, antialiasing, finite spacing, RGB displacement and source variation within a cell all introduce error. The renderer seeks a convincing perceptual approximation rather than lossless coding.
\end{warningbox}

\subsection{Sampling and antialiasing}

Every ideal geometric mask contains sharp boundaries. Raster output samples those boundaries on a finite pixel grid. Without antialiasing, an ideal binary mask

\begin{equation}
M_q^{(0)}(x,y)\in\{0,1\}
\end{equation}

would produce stair-stepping and temporal flicker. A smoothed mask can be modelled as

\begin{equation}
M_q(x,y)=
\left(h_{\mathrm{aa}}*M_q^{(0)}\right)(x,y),
\label{eq:aa}
\end{equation}

where $h_{\mathrm{aa}}$ is a compact reconstruction kernel and $*$ denotes convolution. Supersampling followed by area downsampling is an equivalent practical method.

Smoothing should remove pixel-grid artefacts without blurring neighbouring RGB structures together. This trade-off is especially important for thin diagonal lines and dense spiral turns.

\clearpage
